\begin{pblock}{The problem}
Intertidal flats are among the most productive and fastest-changing
coastal environments, yet the hardest to survey: too shallow for vessels,
too soft to walk, and \term{never wholly visible at any single instant} —
the tide hides part of the surface in every image ever taken of it.

Their elevation therefore cannot be observed. It has to be \emph{inferred},
by watching how the waterline migrates across many images taken at different
tide states, and every method rests on knowing the water level at the moment
each image was acquired.

\vspace{3mm}
\begin{pcallout}
This poster presents \term{MAREA}: the water level is no longer assumed
--- the archive itself measures the interior tide (a distributed tide
gauge made of pixels), corrects the elevation inversion with it, and
delivers the first intertidal DEM with calibrated per-pixel uncertainty,
declared hydraulics, and proven limits.
\end{pcallout}
\end{pblock}
